% =============================================================================
%  EMPIRICAL SECTIONS DRAFT — Sections 4, 5, and 6
%  To be integrated into main.tex (replaces placeholder Sections 4–6)
%  Aligned with notation and style of main.tex
%  Last updated: March 2026
% =============================================================================

% =============================================================
% Section 4 — Empirical Strategy
% =============================================================

\section{Empirical Strategy}
\label{sec:estimation}

This section describes the empirical design that identifies the structural parameters of the transition system developed in Section~\ref{sec:trilemma}. The estimation serves a dual purpose: it provides direct evidence on the transmission channels of pandemic fiscal policy, and it calibrates the parameters ($\alpha_S$, $\psi$, $\eta$, $\alpha_F^{DI}$, $\kappa_F^{CP}$, $\kappa_F^{DI}$, $\gamma_y$) that enter the iLQR optimization in Section~\ref{sec:results}.

% -----------------------------------------------------------------
\subsection{Data and Measurement}
\label{subsec:data}

The estimation uses a balanced panel of 38 OECD member states observed at quarterly frequency from Q1.2020 through Q1.2022---the period during which fiscal deployment was active and vaccination had not yet dissolved the trilemma.

\paragraph{State variables.} The output gap $y_k$ is measured as the percentage deviation of real GDP from its Hodrick-Prescott-filtered trend estimated on the 2015--2019 baseline. Public debt $b_k$ is measured as the first difference of the real debt-to-GDP ratio, normalized by 2019 GDP to isolate changes in the numerator from denominator effects.\footnote{Normalizing by pre-pandemic GDP (2019) avoids the mechanical inflation of the debt ratio that a recession produces through the denominator. This choice follows the specification in equation~(\ref{eq:debt_transition}).} Infection prevalence $\theta_k$ is constructed as the quarterly mean of weekly infection rates, estimated from excess mortality data using country-specific infection fatality rates. Excess mortality $d_k$ follows the WHO methodology \citep{Msemburi2023WHO}.

\paragraph{Control variables.} Containment intensity $S_k$ is the Oxford COVID-19 Government Response Tracker Stringency Index \citep{Hale2021}, averaged within each quarter and population-weighted across subnational jurisdictions where available. The vaccination rate (cumulative doses per capita) controls for the dissolution of the trilemma in later quarters.

\paragraph{Fiscal instruments.} The fiscal variables $F_k^{CP}$, $F_k^{DI}$, and $F_k^{H}$ are constructed from a micro-level database of 2,301 individual fiscal measures classified by transmission channel following the protocol described in Section~\ref{sec:data} and documented in \citet{Codebook}. Each measure is assigned to one of three categories---Capacity Preservation (CP), Demand Injection (DI), or Health (H)---based on its primary economic transmission mechanism, not its recipient. Tax deferrals (IMF Category~3) are excluded from the estimation sample, as their realized fiscal cost approaches zero for solvent taxpayers.\footnote{154 deferral measures are excluded. Their inclusion in a separate robustness specification does not alter the results (Appendix~Table~\ref{tab:rob_deferrals}).} For each country-quarter, the fiscal variables aggregate the measure-level fiscal costs (as \% of 2019 GDP) within each transmission channel. Table~\ref{tab:descriptives} reports summary statistics.

%\input{tables/descriptives.tex}

% -----------------------------------------------------------------
\subsection{Identification}
\label{subsec:identification}

The identification strategy rests on three pillars.

\paragraph{Cross-country variation in fiscal composition.} Countries facing comparable pandemic conditions chose markedly different fiscal mixes. The cross-country correlation between mean stringency and total fiscal deployment is $r = 0.07$---effectively orthogonal. This near-zero correlation implies that fiscal composition was driven by institutional capacity, political preferences, and pre-existing policy frameworks rather than by the contemporaneous severity of the pandemic. It is this variation in the \textit{composition} of the fiscal response, conditional on comparable containment trajectories, that identifies the instrument-specific parameters.

\paragraph{The stringency interaction.} The theoretical model predicts that CP operates through the persistence channel ($\psi S_k y_{k-1}$) and the cushioning channel ($\delta S_k F_k^{CP}$)---both of which require active containment. The empirical counterpart is the interaction $S_k \cdot F_k^{CP}$, which provides within-country, within-quarter variation that is not confounded by common time shocks (absorbed by quarter fixed effects) or time-invariant country characteristics (absorbed by country fixed effects). The interaction identifies $\eta$ and $\delta$ from the joint variation in stringency and fiscal deployment across countries and quarters.

\paragraph{Timing structure.} The theoretical model implies that CP is contemporaneous (firms require support while lockdowns are active) and DI operates with a lag (transfers take time to reach spending). This timing prediction is testable rather than assumed: the lag structure is determined empirically by comparing specifications with $F_{k}^{DI}$, $F_{k-1}^{DI}$, and $F_{k-2}^{DI}$, with the data selecting lag~2 as the specification with the strongest fit.

% -----------------------------------------------------------------
\subsection{Estimating Equations}
\label{subsec:equations}

The estimating equations map directly onto the transition system in Section~\ref{sec:trilemma}. The output equation corresponds to a linearized version of~(\ref{eq:output_full}):
%
\begin{equation}
    y_{ik} = \underbrace{\alpha_S}_{\text{direct cost}} S_{ik}
    + \underbrace{\psi}_{\text{hysteresis}} S_{ik} \cdot y_{i,k-1}
    + \underbrace{\alpha_F^{CP}}_{\text{CP level}} F_{ik}^{CP}
    + \underbrace{\eta}_{\text{cushioning}} S_{ik} \cdot F_{ik}^{CP}
    + \underbrace{\alpha_F^{DI}}_{\text{DI multiplier}} F_{i,k-2}^{DI}
    + \gamma_i + \delta_k + \varepsilon_{ik}
    \label{eq:output_est}
\end{equation}
%
where $\gamma_i$ and $\delta_k$ denote country and quarter fixed effects, respectively. Two-way fixed effects (TWFE) are appropriate for the output equation because $S_k$ is endogenous to the pandemic cycle---quarter fixed effects absorb the common component of both the pandemic shock and the policy response that is not captured by $\theta_k$ and $S_k$ directly.

The debt equation corresponds to~(\ref{eq:debt_transition}):
%
\begin{equation}
    \Delta b_{ik}^R = \underbrace{-\gamma_y}_{\text{auto. stab.}} y_{ik}
    + \underbrace{\kappa_F^{CP}}_{\text{CP cost}} F_{ik}^{CP}
    + \underbrace{\kappa_F^{DI}}_{\text{DI cost}} F_{i,k-1}^{DI}
    + \phi \cdot t + \mu_i + u_{ik}
    \label{eq:debt_est}
\end{equation}
%
where $\mu_i$ denotes country fixed effects and $\phi \cdot t$ a linear time trend. Quarter fixed effects are deliberately excluded from the debt equation: the Mundlak (1978) auxiliary regression rejects the correlated random effects assumption ($p = 0.006$), confirming that country fixed effects are required; however, quarter fixed effects would absorb the common fiscal deployment cycle that constitutes the identifying variation for $\kappa_F^{CP}$ and $\kappa_F^{DI}$ \citep{chetty2020, deb2021}. A linear trend captures the systematic debt dynamics without absorbing the policy variation.

All standard errors are clustered by country (HC1), which is robust to arbitrary within-country serial correlation and heteroskedasticity.


% =============================================================
% Section 5 — Results
% =============================================================

\section{Results}
\label{sec:results_empirical}

% -----------------------------------------------------------------
\subsection{Output Equation}
\label{subsec:output_results}

Table~\ref{tab:output_equation} reports the output equation estimates, building from pooled OLS (Column~1) through country fixed effects (Column~2) to the main two-way fixed effects specification (Column~3). Column~4 replaces $F_{k-2}^{DI}$ with $F_{k-1}^{DI}$ to verify the lag structure.

The main specification (Column~3) yields three findings that map directly onto the theoretical predictions of Section~\ref{sec:trilemma}. First, the interaction $S_k \cdot F_k^{CP}$ is negative and significant ($-0.006$, $p < 0.01$), identifying the combined cushioning-persistence channel $\eta$ from equation~(\ref{eq:rho_eff}). At the sample mean stringency $\bar{S} = 42$, the net marginal effect of CP on output is $\hat{\alpha}_F^{CP} + \hat{\eta} \cdot \bar{S} = 0.260 - 0.006 \times 42 \approx 0.01$---effectively zero. CP is operative precisely during active containment and neutral otherwise, confirming its theoretical role as a lockdown-specific instrument rather than a general-purpose stimulus.

Second, demand injection is significant only at lag~2 ($\hat{\alpha}_F^{DI} = 0.239$, $p < 0.05$) and insignificant at lag~1 ($-0.017$, Column~4). The two-quarter delay is consistent with the implementation chain of transfer programs: authorization, disbursement, and spending each require approximately one quarter, placing the output effect at $k+2$ relative to the deployment quarter.\footnote{This timing is consistent with \citet{chetty2020}, who document that the first-round Economic Impact Payments under the CARES Act took approximately 3--4 weeks to reach households and an additional 4--8 weeks to translate into measurable spending increases.}

Third, the interaction $S_k \cdot y_{k-1}$ is positive and significant ($\hat{\psi} = 0.004$, $p < 0.001$), identifying the lockdown-induced hysteresis parameter from equation~(\ref{eq:rho_eff}). Containment increases the persistence of the output gap: at $\bar{S} = 42$, the effective persistence rises from $\hat{\rho}_y = 0.063$ to $\hat{\rho}_y^{\text{eff}} = 0.063 + 0.004 \times 42 = 0.23$---still below unity but substantially elevated relative to the baseline.

The critical stringency threshold at which CP's marginal output benefit turns negative is $S^* = \hat{\alpha}_F^{CP} / |\hat{\eta}| = 0.260 / 0.006 = 42$, which coincides with the sample mean. This implies that the average OECD economy during the pandemic operated precisely at the boundary of CP effectiveness---additional CP was effective in high-stringency quarters but not in low-stringency ones. This prediction is confirmed by the subsample analysis in Section~\ref{subsec:robustness}.

%\input{tables/output_equation.tex}

% -----------------------------------------------------------------
\subsection{Debt Equation}
\label{subsec:debt_results}

Table~\ref{tab:debt_equation} reports the debt equation estimates. Column~2 is the main specification with pooled CP; Column~3 splits CP into above-the-line (direct grants, wage subsidies) and below-the-line (loans and guarantees, adjusted to the estimated 35\% take-up rate).

The pooled estimate yields $\hat{\kappa}_F^{CP} = 0.191$ ($p < 0.01$): each percentage point of GDP deployed as CP increases the real debt-to-GDP ratio by 0.19~pp per quarter. The automatic stabilizer coefficient is $\hat{\gamma}_y = 0.207$ ($p < 0.001$), implying that a 1~pp widening of the output gap automatically adds 0.21~pp to the quarterly debt change through reduced revenues and increased mandatory transfers.

The split specification (Column~3) reveals that the pooled estimate masks substantial heterogeneity across instrument types. Above-the-line CP (direct grants and wage subsidies, such as Kurzarbeit) carries a debt coefficient of $\hat{\kappa}_{above}^{CP} = 0.464$ ($p < 0.10$), while below-the-line CP (loans and guarantees) carries $\hat{\kappa}_{below}^{CP} = 0.464$ ($p < 0.01$). These near-identical point estimates are misleading, however, because they pool two structurally distinct instruments within below-the-line CP.

%\input{tables/debt_equation.tex}

% -----------------------------------------------------------------
\subsection{Instrument Decomposition: Loans versus Guarantees}
\label{subsec:decomposition}

The sub-component analysis in Table~\ref{tab:decomposition} provides the most novel empirical finding of this paper. Below-the-line CP is disaggregated into preferential loans (Codes~40--41: actual disbursements creating government financial assets) and credit guarantees (Code~43: contingent liabilities that create fiscal costs only when called). This decomposition is applied to both the output and the debt equation.

\paragraph{Output equation (Panel~A).} When the three CP sub-components are each interacted with stringency, the results reveal a sharp separation. Guarantees drive the capacity preservation channel: the interaction $S_k \cdot F_k^{\text{guar}}$ is significant at the 1\% level ($-0.025$, $p = 0.002$), with a level effect of $0.981$ ($p = 0.012$). Loans show an immediate but transitory level effect ($0.253$, $p = 0.003$ at $h=0$ in the LP specification) but no significant stringency interaction. Above-the-line CP is individually insignificant when separated from the other components.

The interpretation is structural: guarantees operate through the announcement channel---the existence of the guarantee framework prevents insolvency cascades and preserves firm-worker matches, regardless of whether individual guarantees are drawn. This is the ``whatever it takes'' mechanism \citep{draghi2012}. The effect is conditional on active containment ($S_k > 0$), confirming that guarantees are a lockdown-specific capacity preservation tool. Loans, by contrast, operate as direct liquidity injections---they stabilize output on impact but do not generate the sustained capacity preservation effect that the stringency interaction identifies.

\paragraph{Debt equation (Panel~B).} The mirror image emerges in the debt equation. Loans are highly significant ($\hat{\kappa}_{\text{loans}} = 0.339$, $p < 0.001$): actual disbursements directly increase government debt. Guarantees are insignificant ($\hat{\kappa}_{\text{guar}} = 0.065$, $p = 0.48$): most guarantees were never called and therefore generated minimal fiscal cost.

\paragraph{Combined finding.} The same below-the-line instrument has opposite signatures across equations:
%
\begin{center}
\begin{tabular}{lcc}
\toprule
 & Output effect & Debt cost \\
\midrule
Guarantees & Significant (via $S \cdot F^{\text{guar}}$) & Insignificant \\
Loans & Insignificant (via $S$ interaction) & Significant ($0.339^{***}$) \\
\bottomrule
\end{tabular}
\end{center}
%
Guarantees preserve output through the announcement channel but create no debt; loans create debt through actual disbursement but have no capacity preservation effect. This decomposition resolves the apparent puzzle in the aggregate data: below-the-line CP appeared simultaneously effective and cheap because it comprised two instruments with fundamentally different transmission mechanisms.

%\input{tables/decomposition.tex}

% -----------------------------------------------------------------
\subsection{Dynamic Effects: Local Projections}
\label{subsec:lp_results}

To trace the full transmission profile of fiscal instruments, I estimate panel Local Projections \citep{jorda2005} at horizons $h = 0, \ldots, 5$. The dependent variable is the cumulative change $y_{i,t+h} - y_{i,t-1}$ for output and $b_{i,t+h}^R - b_{i,t-1}^R$ for debt. The LP sample extends to Q4.2022 to provide response horizons for late-pandemic shocks, with pre-pandemic quarters (2019) anchoring the counterfactual at $F = 0$.

\paragraph{Output dynamics.} The joint LP confirms that the unconditional output response to pooled CP and DI is not statistically significant under TWFE---consistent with the static panel finding that the operative channel is the stringency interaction, not the level. The state-dependent LP, which includes both $F_k^{CP}$ and $S_k \cdot F_k^{CP}$ as regressors, yields significant results at $h = 0$: the CP level effect is $0.241$ ($p = 0.006$) and the interaction is $-0.006$ ($p = 0.022$). Both effects are contemporaneous and decay rapidly, confirming that capacity preservation is a lockdown-period instrument with no lagged dynamics.

As a robustness check, I re-estimate the LP with country fixed effects only (no quarter FE), which preserves the temporal variation absorbed by TWFE. The stringency interaction strengthens to $-0.025$ ($p < 0.001$) and remains significant through $h = 3$ ($-0.006$, $p = 0.009$). The 4$\times$ increase in magnitude between TWFE and country-FE-only specifications bounds the true effect: TWFE provides a conservative lower bound by absorbing both confounders and signal, while country-FE provides an informative upper bound that may include residual common-shock contamination. The qualitative conclusion---CP operates through the stringency interaction---is robust to both specifications.

\paragraph{Debt dynamics (Figure~\ref{fig:lp_main}).} The debt LP produces the strongest dynamic results. Each 1~pp GDP of CP creates $0.16$~pp of debt at $h = 0$, rising monotonically to $0.63$~pp at $h = 5$---with no evidence of mean reversion. The debt effect of CP is persistent and accumulating. DI debt builds more slowly (insignificant until $h = 2$--$3$, reaching $1.00$~pp at $h = 5$), reflecting the implementation lag documented in the static panel.

The CP sub-component LP reveals the dynamic signature of the guarantee channel in the debt equation. Above-the-line CP (grants, wage subsidies) is the costliest instrument, reaching $1.53$~pp debt at $h = 5$ ($p < 0.001$). Loans create immediate and persistent debt ($0.35$ at $h = 0$, $0.68$ at $h = 5$, both $p < 0.01$). Guarantees create debt \textit{only with a 3-quarter delay}: insignificant at $h = 0$--$2$, significant from $h = 3$ ($0.60$, $p = 0.03$) and rising to $1.00$ at $h = 5$ ($p < 0.01$). This delayed debt signature is consistent with the typical timeline of guarantee calling: borrowers default on guaranteed loans 2--4 quarters after the initial disbursement, triggering government payouts.

%\begin{figure}[htbp]
%    \centering
%    \includegraphics[width=\textwidth]{figures/lp_main_2x2.pdf}
%    \caption{Local Projection Impulse Response Functions: Fiscal Policy $\to$ Output Gap (top) and Government Debt (bottom). Shaded bands represent 90\% (dark) and 95\% (light) confidence intervals. Cumulative response per 1~pp GDP fiscal shock. Country-clustered standard errors.}
%    \label{fig:lp_main}
%\end{figure}

\paragraph{Cumulative multipliers.} The LP coefficients at each horizon directly measure the cumulative fiscal multiplier---the total output change per unit fiscal shock. At $h = 5$, the cumulative CP multiplier is effectively zero under TWFE (the output effect decays) while the cumulative DI multiplier is $0.13$ (insignificant). The debt multipliers, by contrast, accumulate: CP reaches $0.63$ and DI reaches $1.00$ at $h = 5$. The implied output-per-debt efficiency ratio declines over time for both instruments, reflecting the permanent debt accumulation without corresponding permanent output gains---a dynamic manifestation of the trilemma.


% -----------------------------------------------------------------
\subsection{Robustness}
\label{subsec:robustness}

The main findings are subjected to extensive robustness checks, reported in full in Appendix~\ref{app:robustness}. I summarize the key results here.

\paragraph{Subsample stability: the stringency split.} The most informative robustness check splits the sample by median pandemic stringency. In high-stringency countries, the $S_k \cdot F_k^{CP}$ interaction is significant ($-0.005$, $p = 0.047$) and CP remains the dominant fiscal channel. In low-stringency countries, CP is completely insignificant ($-0.023$, $S \cdot F^{CP} = +0.001$). This is the cleanest confirmation of the theory's central prediction: capacity preservation is a lockdown-specific policy tool that has no output effect when containment is mild.

\paragraph{Social safety nets.} Splitting by welfare state type (Esping-Andersen typology), DI is insignificant in countries with strong pre-existing safety nets (Nordic and Continental Europe: $-0.068$, $p = 0.84$) but marginally significant in countries with weaker automatic stabilizers ($0.254$, $p = 0.05$). CP is significant in both subsamples ($0.333$ and $0.262$, both $p < 0.06$). This confirms that CP operates through a distinct channel: it preserves productive structures on the firm side, independently of the household-side transfer infrastructure. Discretionary DI substitutes for missing automatic stabilizers; where the safety net is already comprehensive, discretionary DI adds no incremental effect.

\paragraph{Income and fiscal space.} The CP interaction is equally significant in advanced and emerging economies ($-0.007$ and $-0.006$, both $p < 0.02$), suggesting that the capacity preservation mechanism is universal and does not depend on development level. The debt split reveals that high-debt countries relied more on CP ($0.286^{***}$) while low-debt countries could afford DI ($0.284^{*}$)---consistent with the fiscal space hypothesis that CP's lower per-unit cost made it the preferred instrument under constrained fiscal space.

\paragraph{Standard errors and outliers.} HC3 standard errors, cluster bootstrap ($R = 999$), and Driscoll-Kraay corrections yield qualitatively identical inference. Excluding Ireland (``Leprechaun Economics''), Turkey (idiosyncratic inflation dynamics), or both does not alter any coefficient by more than 15\% or change any significance level.

\paragraph{Fixed effects specification.} The LP comparison between TWFE and country-FE-only provides bounds on the output estimates. The $S \cdot F^{CP}$ interaction is significant under both specifications, with the TWFE estimate ($-0.006$) providing a lower bound and the country-FE estimate ($-0.025$) an upper bound. The debt equation results are invariant to the FE choice, as the debt specification already uses country FE only.


% =============================================================
% Summary — Structural Parameters for the iLQR
% =============================================================

\subsection{Structural Parameters for the iLQR}
\label{subsec:parameters_ilqr}

Table~\ref{tab:structural_params} collects the estimated structural parameters that enter the iLQR optimization of Section~\ref{sec:results}. The output equation identifies $\hat{\alpha}_S$, $\hat{\psi}$, $\hat{\alpha}_F^{CP}$, $\hat{\eta}$, and $\hat{\alpha}_F^{DI}$. The debt equation identifies $\hat{\gamma}_y$, $\hat{\kappa}_F^{CP}$, and $\hat{\kappa}_F^{DI}$. Health expenditure $\hat{c}_H$ is calibrated from the classified health measures in the database.

\begin{table}[htbp]
\centering
\caption{Estimated Structural Parameters}
\label{tab:structural_params}
\small
\begin{tabular}{llcl}
\toprule
\textbf{Parameter} & \textbf{Description} & \textbf{Estimate} & \textbf{Source} \\
\midrule
$\hat{\alpha}_S$        & Direct output cost of containment        & $-0.028$          & Table~\ref{tab:output_equation}, Col.~3 \\
$\hat{\psi}$            & Lockdown-induced hysteresis               & $0.004$           & Table~\ref{tab:output_equation}, Col.~3 \\
$\hat{\alpha}_F^{CP}$   & CP level effect on output                & $0.260$           & Table~\ref{tab:output_equation}, Col.~3 \\
$\hat{\eta}$            & CP$\times$S interaction (cushioning)      & $-0.006$          & Table~\ref{tab:output_equation}, Col.~3 \\
$\hat{\alpha}_F^{DI}$   & DI multiplier (lag 2)                    & $0.239$           & Table~\ref{tab:output_equation}, Col.~3 \\
$\hat{\gamma}_y$        & Automatic stabilizer semi-elasticity      & $0.207$           & Table~\ref{tab:debt_equation}, Col.~2 \\
$\hat{\kappa}_F^{CP}$   & Fiscal cost of CP (pooled)               & $0.191$           & Table~\ref{tab:debt_equation}, Col.~2 \\
$\hat{\kappa}_F^{DI}$   & Fiscal cost of DI (lag 1)                & $0.406$           & Table~\ref{tab:debt_equation}, Col.~2 \\
$\hat{c}_H$             & Health expenditure per unit $\theta$     & calibrated        & Health measures (Codes~30--34) \\
\bottomrule
\multicolumn{4}{p{13cm}}{\footnotesize \textit{Notes:} Parameters estimated from the panel specifications in Tables~\ref{tab:output_equation}--\ref{tab:debt_equation}. $\hat{\kappa}_F^{CP}$ is the pooled estimate; the decomposition into above-the-line ($0.464$) and below-the-line ($0.464$) is used in the sensitivity analysis. All estimates use cluster-robust HC1 standard errors.}
\end{tabular}
\end{table}

These parameters, together with the epidemiological parameters ($\rho_\theta$, $\phi_S$, $\delta_\theta$) calibrated from the infection and mortality data, fully specify the transition system for the iLQR. The LP impulse responses provide additional dynamic validation: the monotonically accumulating debt IRF ($0.16 \to 0.63$ over five quarters) and the contemporaneous-then-decaying output IRF confirm that the quarterly transition equations capture the relevant dynamics within the iLQR's planning horizon. The next section takes these estimated parameters to the optimization framework and computes the Pareto frontier across the trilemma dimensions.
